\chapter{Conclusão}\label{chp:CONCLUSAO}

Este capítulo encerra o trabalho com um breve resumo do que foi apresentado, considerações sobre o que foi alcançado em relação ao objetivo proposto na Introdução (Capítulo~\ref{chp:INTRODUCAO}), as limitações identificadas e sugestões para trabalhos futuros.

\section{Considerações finais}

Este trabalho partiu da necessidade de facilitar a busca por atas de reuniões acadêmicas do Departamento de Ciência da Computação da UFRRJ, propondo um chatbot de regras capaz de mapear as perguntas de discentes, docentes e do público acadêmico em geral para as atas mais adequadas ao conteúdo de suas buscas. Para isso, o Capítulo~\ref{chp:FUNDAMENTACAO} situou o problema dentro da literatura sobre chatbots e sobre Busca e Recuperação da Informação (BRI), apresentando a taxonomia dos principais tipos de chatbot disponíveis, o comparativo entre essas opções e os requisitos do projeto, e os modelos clássicos de recuperação, com destaque para o modelo probabilístico e o algoritmo BM25, utilizado pelo Elasticsearch para o ranqueamento padrão de resultados.

A partir dessa fundamentação, o Capítulo~\ref{chp:PROPOSTA} caracterizou o problema central deste trabalho: a desconexão entre o vocabulário livre empregado pelo usuário em uma conversa e a forma como o conteúdo das atas está organizado no índice do Elasticsearch. A solução proposta foi um roteiro guiado de diálogo, que conduz o usuário por um conjunto fechado de critérios de busca, traduz suas respostas em parâmetros estruturados e delimita, de forma explícita, o que pode ou não ser buscado, fechando essa lacuna sem exigir que o usuário conheça a organização do índice ou a sintaxe de consulta do motor de busca.

O Capítulo~\ref{chp:EXPERIMENTOS} descreveu a implementação dessa proposta, composta pelo Botpress, responsável pelo fluxo de diálogo; pelo indexador, desenvolvido em Node.js com Express, responsável por extrair o conteúdo das atas em PDF e montar a consulta ao Elasticsearch; pelo Elasticsearch propriamente dito, responsável pelo armazenamento e pelo ranqueamento via BM25; e por um bot do Discord, oferecido como canal alternativo de interação. Foi detalhado o mapeamento de cada um dos sete critérios de busca do roteiro (departamento, disciplina, docentes, discentes, período, processo/edital e progressão), além do campo de assunto geral, para cláusulas da consulta em \textit{Query DSL}, com pesos distintos conforme a especificidade de cada critério, e a aplicação de um cálculo adicional de proporcionalidade para reordenar os três melhores resultados antes de sua exibição ao usuário.

Em relação ao objetivo definido na Introdução — disponibilizar uma ferramenta conversacional capaz de interpretar perguntas em linguagem natural e mapeá-las para as atas institucionais mais pertinentes —, o trabalho alcançou esse resultado por meio da combinação de um chatbot baseado em regras com um roteiro guiado de diálogo, com as definições dos paramêtros que seriam oferecidos para o usuário realizar a personalização da busca. O comparativo do Capítulo~\ref{chp:FUNDAMENTACAO} (Tabela~\ref{tab:comparativo-chatbots}) evidenciou que a escolha por um chatbot baseado em regras é coerente com o escopo fechado do problema, com a ausência de orçamento para licenciamento de ferramentas e com a necessidade de integração via API com o Elasticsearch. O roteiro guiado de diálogo, por sua vez, mostrou-se uma solução viável para o problema de geração de consultas a partir de uma interação conversacional, ao definir de forma explícita quais critérios podem ser buscados e como cada resposta do usuário se traduz em um componente da consulta, sem exigir do usuário qualquer conhecimento sobre a organização do índice ou sobre o BM25.

\section{Limitações e trabalhos futuros}

Por se tratar de um chatbot baseado em regras, o sistema opera sobre um vocabulário e um conjunto de critérios de busca fechados e definidos previamente no roteiro de diálogo. Conforme discutido nos Capítulos~\ref{chp:FUNDAMENTACAO} e \ref{chp:PROPOSTA}, essa característica é adequada ao escopo do problema, mas limita a interpretação de perguntas que não se encaixem nos sete critérios previstos (departamento, disciplina, docentes, discentes, período, processo/edital e progressão) ou no campo de assunto geral, o qual concentra qualquer informação fornecida pelo usuário que não corresponda a esses critérios ou a um período de datas reconhecível.

A recuperação textual implementada utiliza exclusivamente o BM25, por meio do Elasticsearch, e depende da correspondência lexical entre os termos informados pelo usuário e o conteúdo das atas, ainda que com tolerância a variações ortográficas em alguns campos (\texttt{fuzziness: AUTO}). A busca por similaridade semântica, baseada em \textit{embeddings}, foi mencionada na Fundamentação como mecanismo complementar ao BM25, situando-se dentro da família de modelos vetoriais discutida no Capítulo~\ref{chp:FUNDAMENTACAO}, mas não foi implementada neste trabalho. Sua incorporação ao indexador, complementando os resultados do BM25, é apontada como uma extensão natural, com potencial para capturar relações de significado entre o vocabulário do usuário e o conteúdo das atas que a correspondência lexical do BM25 não captura, especialmente em casos de sinônimos, termos coloquiais ou descrições indiretas de um assunto.

Outra limitação está relacionada ao próprio conjunto de critérios de busca: ele foi definido a partir de aspectos recorrentes do conteúdo das atas institucionais da UFRRJ, mas novos tipos de informação que venham a ganhar relevância no conteúdo de atas futuras exigiriam a inclusão de uma nova etapa no roteiro de diálogo e do componente de busca correspondente na consulta ao Elasticsearch, conforme discutido no Capítulo~\ref{chp:PROPOSTA}. Embora o roteiro tenha sido projetado para tornar essa extensão simples de implementar, a manutenção e a ampliação do escopo permanecem como trabalho a ser realizado.

Por fim, sugere-se como trabalho futuro a realização de uma avaliação de usabilidade com usuários reais do Departamento de Ciência da Computação da UFRRJ, envolvendo discentes, docentes e o público acadêmico em geral, de modo a verificar, em uso real, a efetividade do roteiro guiado na recuperação das atas desejadas e a identificar critérios de busca ou formulações de pergunta não previstos no escopo atual.
